\section{\termCuspTrajectorySpace{}的离壳描述}
\label{sec:off-shell-cusp-trajectory-spaces}

下面给出尖点流线模空间的离壳描述：
\begin{equation}
\begin{aligned}
 &\widetilde{\mathcal M}\bigl(p,q;B_0;
 \{\alpha_i\},\{\beta_{0,j}\},\{\beta_{1,k}\}\bigr)\\
 &\quad:=\widetilde{\mathcal M}_{m;n_0,n_1}(p,q;B_0)_{\mathrm{ev}}
 \mathop{\times}_{\mathrm{ev}}
 \Bigl(\prod_i\mathcal M_1(J_{(\tau_i,t_i)};\alpha_i)\\
 &\hspace{9em}\times\prod_j\mathcal M_1(L,J_0;\beta_{0,j})
 \times\prod_k\mathcal M_1(L,J_1;\beta_{1,k})\Bigr).
\end{aligned}
\label{eq:ch15-cusp-trajectory-moduli}
\end{equation}
与第~\ref{sec:off-shell-framework-smooth-floer}节一样，我们要找出这个纤维积的离壳对应物。
对第一个因子 $\widetilde{\mathcal M}_{(m;n_0,n_1)}(p,q;B_0)$，考虑乘积
\[
 \mathcal F_{\delta;(m;n_0,n_1)}^{k,p}(x,y;B)
 :=\mathcal F_\delta^{k,p}(x,y;B)\times\Theta^m
 \times(\partial_0\Theta)^{n_0}\times(\partial_1\Theta)^{n_1}.
\]
除\personCauchy{}--\personRiemann{}截面 $\overline\partial_J$ 外，还有取值映射
\[
 \operatorname{ev}_{\mathcal F}:
 \mathcal F_{\delta;(m;n_0,n_1)}^{k,p}(x,y;B)
 \longrightarrow M^m\times L_0^{n_0}\times L_1^{n_1},
\]
定义为
\begin{equation}
\begin{aligned}
 &\operatorname{ev}_{\mathcal F}
 \bigl(u,\{(\tau_i,t_i)\},\{\tau_{0,j}\},\{\tau_{1,k}\}\bigr)\\
 &\qquad=\bigl(\{u(\tau_i,t_i)\},\{u(\tau_{0,j},0)\},
 \{u(\tau_{1,k},1)\}\bigr).
\end{aligned}
\label{eq:off-shell-strip-evaluation}
\end{equation}
在离壳描述的讨论中，我们只考虑光滑映射和光滑截面；相应空间的适当完备化留给读者。

对第二个因子 $\prod_i\mathcal M_1(J_{(\tau_i,t_i)};\alpha_i)$，考虑参数化模空间
\[
 \mathcal M_1(J,\alpha;\mathrm{para})
 :=\coprod_{(s,t)\in[0,1]^2}\{(s,t)\}\times\mathcal M_1(J_{(s,t)},\alpha).
\]
模空间 $\widetilde{\mathcal M}_1(J,\alpha;\mathrm{para})$ 的离壳对应物为
\[
 \mathcal F_1(\alpha;\mathrm{para})
 :=W^{k,p}(\Sigma,M;\alpha)\times\Sigma\times[0,1]^2.
\]
我们有取值映射
\[
 \operatorname{ev}^{S}:\mathcal F_1(\alpha;\mathrm{para})\longrightarrow M\times[0,1]^2,
\]
其定义为
\begin{equation}
 \operatorname{ev}^{S}(v,z,s,t)=(v(z),s,t),
 \label{eq:parameterized-sphere-evaluation}
\end{equation}
以及参数化\personCauchy{}--\personRiemann{}截面
\[
 \overline\partial^{\mathrm{para}}:
 \mathcal F_1(\alpha;\mathrm{para})\longrightarrow\mathcal L,
\]
其定义为
\[
 \overline\partial^{\mathrm{para}}(v,z,s,t)
 =\overline\partial_{J_{(s,t)}}(v)\in W^{k-1,p}(v^*TM).
\]
这里，原书把 $\mathcal F_1^{(k,p)}(\alpha;\mathrm{para})$ 写为
\[
 W^{k,p}(\Sigma\times[0,1]^2,M;\alpha)\times\Sigma,
\]
并把 $C^\infty(\Sigma\times[0,1]^2,M;\alpha)$ 写成乘积
\[
 W^{k,p}(\Sigma,M;\alpha)\times[0,1]^2.
\]
\footnote{译者注：本段两式在映射空间的正则性符号及参数放置上彼此不一致，
也与紧邻其前的 $\mathcal F_1(\alpha;\mathrm{para})$ 定义不一致。译文照录两式，
不擅自统一。}

对~\eqref{eq:ch15-cusp-trajectory-moduli}的第三、第四个因子，考虑模空间
$\widetilde{\mathcal M}_1(L,J_0;\beta)$ 与
$\widetilde{\mathcal M}_1(L,J_1;\beta)$。
\footnote{译者注：原书在此把两者都印作
$\widetilde{\mathcal M}_1(L,J_0;\beta)$；但上下文分别对应 $L_0,J_0$ 与
$L_1,J_1$，同页后续直和也明确使用 $J_0,J_1$。译文据此补正第二项。}
$\widetilde{\mathcal M}_1(L,J_0;\beta)$ 的离壳对应物为
\[
 W^{k,p}\bigl((D^2,\partial D^2),(M,L)\bigr)\times\partial D^2,
\]
而取值映射为
\[
 \operatorname{ev}^{D}:W^{k,p}\bigl((D^2,\partial D^2),(M,L)\bigr)
 \times\partial D^2\longrightarrow L,
 \qquad \operatorname{ev}^{D}(w,\theta)=w(\theta).
\]
合并上述取值映射，得到
\begin{equation}
 \operatorname{ev}:
 \prod_i\mathcal F_1(J,\alpha_i)
 \times\prod_j\mathcal F_1(L_0,J_0;\beta_{0,j})
 \times\prod_k\mathcal F_1(L_1,J_1;\beta_{1,k})
 \longrightarrow M^m\times L_0^{n_0}\times L_1^{n_1},
 \label{eq:combined-bubble-evaluation}
\end{equation}
其中
\[
 \operatorname{ev}=(\prod_i\operatorname{ev}_i^S,
 \prod_j\operatorname{ev}_j^D,\prod_k\operatorname{ev}_k^D).
\]
另一方面，有合并的\personCauchy{}--\personRiemann{}截面
\begin{equation}
\begin{aligned}
 \overline\partial_J:
 &\prod_i\mathcal F_1(J,\alpha_i)
 \times\prod_j\mathcal F_1(L_0,J_0;\beta_{0,j})
 \times\prod_k\mathcal F_1(L_1,J_1;\beta_{1,k})\\
 &\longrightarrow
 \mathcal L(x,y;B)\oplus\bigoplus_i\mathcal L(\alpha_i)
 \oplus\bigoplus_j\mathcal L(L_0,\beta_{0,j})
 \oplus\bigoplus_k\mathcal L(L_1,\beta_{1,k}),
\end{aligned}
\label{eq:combined-cauchy-riemann-section}
\end{equation}
它是直和
\[
 \overline\partial_J=
 \left(\overline\partial_J^1,
 \bigoplus\overline\partial^{\mathrm{para},J},
 \bigoplus_j\overline\partial_{J_0}^2,
 \bigoplus_k\overline\partial_{J_1}^2\right).
\]
于是有恒等式
\begin{equation}
\begin{aligned}
 &\widetilde{\mathcal M}_{(m;n_0,n_1)}
 (x,y;B_0;\{\alpha_i\},\{\beta_{0,j}\},\{\beta_{1,k}\})\\
 &\qquad=\overline\partial_J^{-1}(0)
 \cap(\operatorname{ev}_{(m;n_0,n_1)},\operatorname{ev})^{-1}(\Delta).
\end{aligned}
\label{eq:cusp-space-zero-set-fiber-product}
\end{equation}

\stfhIndex{I192}{Xuqiekongjian}{虚切空间（virtual tangent space）}
现在描述
\[
 \widetilde{\mathcal M}_{(m;n_0,n_1)}
 (x,y;B_0;\{\alpha_i\},\{\beta_{0,j}\},\{\beta_{1,k}\})
\]
的\termVirtualTangentSpace{}。令
\[
 \mathbf u=\bigl(u,\{v_i\},\{w_{0,j}\},\{w_{1,k}\};
 \{(\tau_i,t_i)\},\{(\tau_j,0)\},\{(\tau_k,1)\}\bigr)
\]
属于上述模空间。受其纤维积描述启发，考虑映射
\[
 \bigl(\overline\partial_J,
 (\operatorname{ev}_{(m;n_0,n_1)},\operatorname{ev})\bigr)
\]
在 $\mathbf u$ 处的线性化；它作用在
\[
\begin{aligned}
 &\mathcal F_{\delta;(m;n_0,n_1)}^{k,p}(x,y;B)
 \times\prod_i\mathcal F_1(J,\alpha_i)
 \times\prod_j\mathcal F_1(L_0,J_0;\beta_{0,j})\\
 &\hspace{12em}\times\prod_k\mathcal F_1(L_1,J_1;\beta_{1,k})
\end{aligned}
\]
上。后者在 $\mathbf u$ 处的切空间是直和
\[
\begin{aligned}
 &W_{\delta;(m;n_0,n_1)}^{k,p}
 (u^*TM;(\partial_0u)^*TL_0,(\partial_1u)^*TL_1)\\
 &\quad\oplus\bigl(W_1^{k,p}(v_i^*TM)\bigr)^{\oplus m}
 \oplus\bigl(W_1^{k,p}(w_0^*TM,(\partial w_0)^*TL_0)\bigr)^{\oplus n_0}\\
 &\quad\oplus\bigl(W_1^{k,p}(w_1^*TM,(\partial w_1)^*TL_1)\bigr)^{\oplus n_1}.
\end{aligned}
\]
这里记
\[
\begin{aligned}
 &W_{\delta;(m;n_0,n_1)}^{k,p}
 (u^*TM;(\partial_0u)^*TL_0,(\partial_1u)^*TL_1)\\
 &\quad=W_\delta^{k,p}(u^*TM;(\partial_0u)^*TL_0,(\partial_1u)^*TL_1)
 \oplus\bigl((\mathbb R^2)^m\oplus\mathbb R^{n_0}\oplus\mathbb R^{n_1}\bigr),
\end{aligned}
\]
以及
\[
 W_1^{k,p}(w_0^*TM,(\partial w_0)^*TL_0)
 =W^{k,p}(w_0^*TM,(\partial w_0)^*TL_0)\oplus T\partial D^2,
\]
对 $W_1^{k,p}(w_1^*TM,(\partial w_1)^*TL_1)$ 类似，对
$W_1^{k,p}(v_i^*TM)$ 则更为直接。

于是有线性映射
\[
\begin{aligned}
 E_{\mathbf u}^1:
 &W_{\delta;(m;n_0,n_1)}^{k,p}
 (u^*TM;(\partial_0u)^*TL_0,(\partial_1u)^*TL_1)\\
 &\longrightarrow L_\delta^{k,p}(u^*TM)\oplus
 \left((u^*TM)^{\oplus m}\oplus((\partial_0u)^*TL_0)^{\oplus n_0}
 \oplus((\partial_1u)^*TL_1)^{\oplus n_1}\right),
\end{aligned}
\]
\footnote{译者注：原书把 $E_{\mathbf u}^1$ 的目标第一因子写成
$L_\delta^{k,p}(u^*TM)$；但 $D\overline\partial_J(u)$ 是一阶算子，且
定理~\ref{thm:weighted-off-shell-banach-framework}所定义的目标为 $L_\delta^{k-1,p}$。
译文保留原书指标。}
定义为
\[
\begin{aligned}
 &E_{\mathbf u}^1(\xi,\{a_i\},\{b_{0,j}\},\{b_{1,k}\})\\
 &\quad=\bigl(D\overline\partial_J(u)(\xi),
 \{du(\tau_i,t_i)a_i\},
 \{du(\tau_j,0)b_{0,j}\},
 \{du(\tau_k,1)b_{1,k}\}\bigr).
\end{aligned}
\]
\footnote{译者注：原书最后一项印作 $\{du(\tau_k,0)b_{0,k}\}$，与定义域中的
$b_{1,k}$、边界分支 $\partial_1u$ 及紧接着的匹配条件均冲突；译文补正为
$\{du(\tau_k,1)b_{1,k}\}$。}
还有
\[
\begin{aligned}
 E_{\mathbf u}^2:
 &\bigl(W_1^{k,p}(v_i^*TM)\bigr)^{\oplus m}
 \oplus\bigl(W_1^{k,p}(w_0^*TM,(\partial w_0)^*TL_0)\bigr)^{n_0}\\
 &\oplus\bigl(W_1^{k,p}(w_1^*TM,(\partial w_1)^*TL_1)\bigr)^{n_1}\\
 &\longrightarrow\bigoplus_i\bigl(C^\infty(v_i^*TM)\oplus v_i^*TM\bigr)\\
 &\quad\oplus\bigoplus_j\bigl(W^{k-1,p}(w_{0,j}^*TM)
 \oplus(\partial w_{0,j})^*TL_0\bigr)\\
 &\quad\oplus\bigoplus_k\bigl(W^{k-1,p}(w_{1,k}^*TM)
 \oplus(\partial w_{1,k})^*TL_1\bigr),
\end{aligned}
\]
它由合并算子
\[
\begin{aligned}
 E_{\mathbf u}^2=\bigl(&\{(D\overline\partial_{J_{(\tau,t)}}(v_i),dv_i(z_i))\},\\
 &\{(D\overline\partial_{J_0}(w_{0,j}),dw_{0,j}(\theta_{0,j}))\},\\
 &\{(D\overline\partial_{J_1}(w_{1,k}),dw_{1,k}(\theta_{1,k}))\}\bigr)
\end{aligned}
\]
定义。

回忆 $\mathbf u$ 满足匹配条件
\[
\begin{aligned}
 &\operatorname{ev}_{(m;n_0,n_1)}
 (u,\{(\tau_i,t_i)\},\{(\tau_j,0)\},\{(\tau_k,1)\})\\
 &\qquad=\operatorname{ev}
 (\{(v_i,z_i)\},\{(w_{0,j},\theta_{0,j})\},
 \{(w_{1,k},\theta_{1,k})\}),
\end{aligned}
\]
这等价于
\begin{equation}
 u(\tau_i,t_i)=v_i(z_i),\qquad
 u(\tau_j,0)=w_{0,j}(\theta_j),\qquad
 u(\tau_k,1)=w_{1,k}(\theta_k).
 \label{eq:cusp-component-matching}
\end{equation}
\footnote{译者注：原书公式~(15.2.13) 最后一项写作
$w_{0,k}(\theta_k)$，但相应气泡属于 $L_1$ 端且前后均记作 $w_{1,k}$；译文据此补正。}
特别地，
\[
 T_{u(\tau_i,t_i)}M=T_{v_i(z_i)}M,
 \qquad T_{u(\tau_j,0)}L_0=T_{w_{0,j}(\theta_j)}L_0,
 \qquad T_{u(\tau_k,1)}L_1=T_{w_{1,k}(\theta_k)}L_1.
\]
\footnote{译者注：原书第三个等式左端印作 $T_{u(\tau_k,1)}L_0$，
与 $u(\tau_k,1)\in L_1$ 及右端 $TL_1$ 冲突；译文补正为 $TL_1$。}
因此，可以在上述总切空间中定义由下列条件刻画的子空间：
\begin{equation}
 \begin{aligned}
 du(\tau_i,t_i)(\xi)-dv_i(z_i)&=0,\\
 du(\tau_j,0)(\xi)-dw_{0,j}(\theta_j)&=0,\\
 du(\tau_k,1)(\xi)-dw_{1,k}(\theta_k)&=0.
 \end{aligned}
 \label{eq:linearized-cusp-matching}
\end{equation}

现在可以定义
$\mathcal M_{(m;n_0,n_1)}(x,y;B_0,\{\alpha_i\},\{\beta_{0,j}\},\{\beta_{1,k}\})$
在 $\mathbf u$ 处的\termVirtualTangentSpace{}。

\begin{definition}
\label{def:virtual-tangent-cusp-space}
定义
\[
 \operatorname{Index}E_{\mathbf u}
 =\ker E_{\mathbf u}-\operatorname{coker}E_{\mathbf u}
\]
为上述模空间在 $\mathbf u$ 处的\termVirtualTangentSpace{}。
\end{definition}
